% 10 · t-SNE — neighbour probabilities and the crowding problem
\section{t-SNE: A MAP OF NEIGHBOURHOODS}
\label{sec:tsne}

\deckslides{slides 32--38 --- the two probability models and the descent between them.}

Everything so far has treated the 16-dimensional abundance vector as a space
in which distance means something. That assumption is weaker than it looks. In
high dimensions, distances concentrate: the ratio of the farthest to the
nearest neighbour of a point tends to 1, so ``close'' and ``far'' describe
almost the same thing and a density-based method loses contrast. t-SNE
\citep{vanderMaaten:08}, building on stochastic neighbour embedding
\citep{Hinton:02}, answers this by giving up on preserving distances at all
and preserving \emph{neighbourhoods} instead --- then drawing the result in two
dimensions.

\begin{figure}[htbp]\centering
  \begin{tabular}{@{}c@{\hskip5pt}c@{\hskip5pt}c@{}}
    \fig[0.30]{tsne_step1.png} &
    \fig[0.30]{tsne_step2.png} &
    \fig[0.30]{tsne_step3.png}
  \end{tabular}
  \caption{The three moves of t-SNE, from the deck: \panel{a} in the high-D
  space, convert distances to neighbour probabilities (here the point's
  neighbourhoods are drawn as spheres of decreasing probability); \panel{b} in
  the 2-D map, do the same conversion with a heavy-tailed kernel; \panel{c}
  move the points downhill on the Kullback--Leibler divergence until the two
  sets of probabilities agree as closely as they can.}
  \label{fig:tsnesteps}
\end{figure}

\subsection{From distances to probabilities}

For each point $i$, t-SNE defines the probability that $j$ is its neighbour in
the high-dimensional space as a Gaussian in the distance,

\begin{equation}
p_{j \mid i} = \frac{\exp\!\left(-\|\mathbf{x}_i -
\mathbf{x}_j\|^2 / 2\sigma_i^2\right)}
{\sum_{k \neq i} \exp\!\left(-\|\mathbf{x}_i - \mathbf{x}_k\|^2 /
2\sigma_i^2\right)} ,
\label{eq:phigh}
\end{equation}

and symmetrises, $p_{ij} = (p_{j|i} + p_{i|j})/2n$. The novelty is $\sigma_i$:
each point gets its \emph{own} bandwidth, chosen so that the neighbourhood
distribution has a fixed entropy. The user-facing parameter is that entropy,
expressed as \emph{perplexity},

\begin{equation}
\mathrm{Perp}(P_i) = 2^{H(P_i)}, \qquad
H(P_i) = -\sum_j p_{j|i} \log_2 p_{j|i},
\label{eq:perplexity}
\end{equation}

found by binary search on $\sigma_i$. Perplexity is the effective number of
neighbours each point is required to have, and it is the parameter you must
justify. \citet{vanderMaaten:08} suggest 5--50; a widely used default is
$\sqrt{n}$ or $n/100$; our pipeline's default is 30, and the swap to 50
measurably improved the layout for our $\sim\!5\,000$-star samples.

\begin{marginnote}[]
\entry{Perplexity}{The effective number of neighbours each point is given in
t-SNE's high-dimensional probability distribution
(\textbf{Equation~\ref{eq:perplexity}}). It is the only required
hyperparameter, and it is a statement about the scale of structure you are
looking for.}
\entry{Crowding problem}{In a low-dimensional map there is not enough room to
place moderately distant points at moderate distances. t-SNE's heavy-tailed
low-dimensional kernel fixes it by making the map's tail fat.}
\end{marginnote}

\subsection{The low-dimensional half, and the crowding problem}

The map is a set of 2-D positions $\mathbf{y}_i$. Their neighbour
probabilities use a Student-$t$ kernel with one degree of freedom,

\begin{equation}
q_{ij} = \frac{\left(1 + \|\mathbf{y}_i - \mathbf{y}_j\|^2\right)^{-1}}
{\sum_{k \neq l} \left(1 + \|\mathbf{y}_k - \mathbf{y}_l\|^2\right)^{-1}} ,
\label{eq:qlow}
\end{equation}

and the layout is found by minimising the Kullback--Leibler divergence
$D_{\mathrm{KL}}(P \,\|\, Q)$ by gradient descent. The heavy tail in
\textbf{Equation~\ref{eq:qlow}} is the whole point. With a Gaussian map
kernel, the volume available at radius $r$ in two dimensions grows only as
$r$, while the volume of a shell in 16 dimensions grows as $r^{15}$: a point
that should be moderately distant from its neighbours has nowhere to go, and
the optimiser either piles dissimilar points on top of each other or flings
them arbitrarily far away. This is the crowding problem, and the $t_1$ kernel
solves it by allowing large distances to be cheap: $q_{ij} \sim
1/\|\mathbf{y}_i - \mathbf{y}_j\|^2$, so a point can sit far from another at
almost no cost in divergence. The resulting gradient has an intuitive form ---
each point is pulled towards its high-D neighbours and pushed away from the
others, with the repulsion weak at long range --- and it is what
\textbf{Figure~\ref{fig:tsnefilm}} animates.

\begin{figure}[htbp]\centering
  \fig[0.94]{tsne_film.png}
  \caption{The descent, four frames of the deck's animation: from the initial
  (PCA) layout, through the early exaggeration phase in which clusters
  coalesce while empty space opens up between them, to the settled map. The
  exaggeration phase is a device, not physics: it deliberately over-states
  attraction at the start to escape a poor initial arrangement.}
  \label{fig:tsnefilm}
\end{figure}

\begin{issues}[HOW TO READ A t-SNE MAP]
The picture is a map of neighbour probabilities, and the following are all
\emph{not} readable from it. (i) Distances between clusters: only the relative
order of neighbourhoods is preserved, so a gap can appear between two groups
whose members are mutual neighbours. (ii) Cluster sizes and densities: t-SNE
normalises local density away, so a tight blob and a loose one can have the
same number of members. (iii) Global geometry: rotation, reflection and the
arrangement of clusters are arbitrary. (iv) The absence of a gap: a method
that must place every point somewhere will place well-separated points on the
same island if the perplexity is wrong.
\end{issues}

\subsection{What t-SNE does not do}

It does not cluster. The output is a 2-D embedding, and the visual groupings
are an artefact of a continuous optimisation. The workflow that has become
standard --- and that our pipeline uses --- is to treat the embedding as a
feature space and cluster \emph{that} with a density-based method, normally
HDBSCAN$^*$ (\S\ref{sec:hdbscan}) or EVoC (\S\ref{sec:evoc}). The
distinction matters when quoting results: a t-SNE row in our tables is
really ``t-SNE followed by HDBSCAN$^*$'', and both steps contribute to the
score.

Three further practical facts, each of which we have hit:

\begin{itemize}
\item \textbf{Initialisation changes the answer.} The original implementation
  started from a random layout; the modern default starts from a PCA of the
  input. Both are stochastic in practice, and on our abundance matrix the
  difference is not subtle: see the row-order audit of \S\ref{sec:validation},
  where the same data gave homogeneity between 0.218 and 0.560 depending on
  the order of the input rows.
\item \textbf{The distance metric is fixed to Euclidean.} t-SNE has no
  metric parameter, which is why \S\ref{sec:data}'s decision to make
  Euclidean distance and cosine similarity equivalent on our matrix is not a
  detail: it is the only reason the t-SNE, UMAP and EVoC arms are comparable
  at all.
\item \textbf{Cost grows faster than linearly.} The original $O(n^2)$
  formulation is unusable past a few thousand stars; the Barnes--Hut
  approximation \citep{vanderMaaten:14} brings it to
  $O(n \log n)$ per iteration, at the cost of the order dependence noted
  above. Our field runs embed 30\,107 stars in this regime.
\end{itemize}

\subsection{t-SNE on real clusters}

The published anchor for this section is \citet{Kos:17}, who applied t-SNE to
GALAH abundances --- 13 elements, calibrated with The Cannon --- and recovered 7
of the 9 clusters in their sample, the Pleiades among them, as coherent groups in
a two-dimensional map. \textbf{Figure~\ref{fig:kostse}} reproduces their figure,
and it is worth being precise about what it shows: a known cluster, its members
in red, sitting inside a map built from the surrounding field. It is a positive
result, and it is also a reminder of the asymmetry in this literature: the
clusters were selected because they were already known, and the method is judged
on whether it finds them again.

Our own cluster-only results for the abundance matrix illustrate the failure
mode the caveats above predict, and they do it in two registers. The single run
of \textbf{Table~\ref{tab:baseline}} gives the t-SNE arm the highest completeness
of the three methods ($0.879$) and the lowest homogeneity ($0.292$), the
signature of a map in which most clusters have been drawn into one region. The
seven-seed average of \textbf{Table~\ref{tab:honest}} looks better ($0.560$
homogeneity, $0.601$ completeness) --- and \S\ref{sec:validation} shows that the
same matrix moves between $0.218$ and $0.560$ on homogeneity, and between 2 and
22 groups, when only the order of the input rows changes. The neighbourhoods are
preserved, and the neighbourhoods were not separated to begin with.

\begin{figure}[htbp]\centering
  \fig[0.62]{kos_pleiades_tsne.png}
  \caption{t-SNE of GALAH abundances around the Pleiades, as published by
  \citet{Kos:17} and used in the lecture: the red points are stars known to be
  cluster members, and the inset zooms in on where they concentrate, marking the
  two groups the method finds (A, B) and the two member candidates it adds
  (1, 2). A known cluster, recovered from field data --- the benchmark that
  strong chemical tagging has to beat.}
  \label{fig:kostse}
\end{figure}

\begin{issues}[EXERCISES]
\begin{enumerate}
\item Implement the binary search for $\sigma_i$ in
  \textbf{Equation~\ref{eq:perplexity}} that fixes perplexity at a target
  value for every point. Plot the resulting $\sigma_i$ against local density
  on the DR19 matrix. Does the adaptive bandwidth do what it is supposed to
  do?
\item Embed the member matrix with perplexities 5, 30 and 100, and for each
  compute the kNN purity of \S\ref{sec:knn} plus the cluster-only homogeneity.
  Which perplexity would you choose, on which evidence, and does the choice
  change the paper's conclusion?
\item Take the t-SNE embedding of the abundances and run HDBSCAN$^*$ with
  \texttt{min\_cluster\_size} = 5, 15 and 50. Report the number of groups and
  the largest group fraction each time. Identify which run is degenerate and
  say what the degeneracy looks like in the map.
\item The low-dimensional kernel is $q_{ij} \propto (1 +
  \|\mathbf{y}_i-\mathbf{y}_j\|^2)^{-1}$. Replace it with a Student-$t$ kernel
  of $\nu = 5$ degrees of freedom, re-derive the gradient, and explain what
  you expect to change in the map's long-range behaviour.
\end{enumerate}
\end{issues}
